\section{Pengantar Proposisi dan Pemodelan Kalimat Deklaratif}

Sepanjang sejarah peradaban, bahasa menjadi jembatan utama manusia mentransmisikan ide, fakta, emosi, hingga perintah. Namun di mata seorang matematikawan atau perancang mesin komputasi (*computer scientist*), tidak semua tumpahan kosakata manusia memiliki bobot nilai fungsional yang bisa diolah \cite{rosen2019}. Mesin tidak mengerti simpati, kemarahan, atau ambiguitas; ia mutlak hanya bereaksi pada \textbf{kebenaran (True)} dan \textbf{kesalahan (False)}.

Untuk mendikte aliran data kepada mesin biner, kita harus menyaring bahasa manusia ke dalam unit-unit pernyataan terkecil yang tegas nasib nilainya. Unit elementer inilah yang disebut sebagai \textbf{Proposisi}.

\subsection{Definisi Proposisi}

\begin{definisi}[Proposisi Logika]
\logic{Proposisi} (sering diistilahkan sebagai \textbf{Kalimat Tertutup}) adalah suatu wujud \textbf{kalimat deklaratif} (pernyataan yang menerangkan atau mengklaim sebuah fakta) yang \textbf{hanya memiliki KEPATENAN SATU NILAI}: Entah Mutlak Benar Saja (True), atau Mutlak Salah Saja (False), namun haram dan mustahil mengemban keduanya sekaligus pada ruang waktu dimensi yang sama.
\end{definisi}

Konsep ketegasan nilai mutlak tunggal ini dalam Aljabar meminjam \textit{Hukum Tiada Jalan Tengah (Law of Excluded Middle)}. Entitas proposisi menolak keras segala bentuk area abu-abu, ketidakpastian opini, hingga subjektivitas sentimen.

\begin{contoh}[Kalimat Sah Proposisi]
Mari kita uji lima klaim deklaratif berikut:
\begin{enumerate}
    \item "Menara Eiffel menjulang megah menancap tepat di pusat Ibukota Jepang." $\rightarrow$ \textit{(Ini sah sebagai Proposisi, meski dia mengandung fakta pemalsuan letak geografis. Hal terpenting, ia \textbf{bisa dibuktikan mutlak SALAH.})}
    \item "Tiga ditambahkan dua mengkalkulasi hasil final Lima ($3 + 2 = 5$)." $\rightarrow$ \textit{(Proposisi yang absolut BENAR berdasar postulat aritmatika.)}
    \item "Kucing bernapas menghirup rute biologis paru-paru mamalia." $\rightarrow$ \textit{(Proposisi absolut BENAR berdasar observasi sains botani/zoologi.)}
    \item "Jumlah populasi gajah Sumatra di titik konservasi hari ini menembus persis $14.238$ ekor." $\rightarrow$ \textit{(Ini tetaplah proposisi empiris sah! Meski Anda pribadi ragu dan tidak tahu keakuratannya di lapangan detik ini, namun esensinya nilai kebenaran mutlaknya SUNGGUH ADA (entah T atau F) di sisi kenyataan semesta sesungguhnya.)}
\end{enumerate}
\end{contoh}

\subsection{Kategori \textit{Non-Proposisi} (Kalimat Sampah Logika)}

Bila syarat utama proposisi menuntut sifat **deklaratif pasti**, maka di luar batas kerajaan Logika Proposisional, berserakanlah varian format bahasa manusia yang ditolak masuk diproses oleh unit *compiler* kita \cite{wiki_first_order_logic}.

\begin{contoh}[Kalimat yang Mustahil Dikalkulasi]
Berikut ragam bahasa yang tidak diakui di pengadilan matematika:
\begin{enumerate}
    \item \textbf{Kalimat Interogatif (Tanya):} "Berapakah suhu panas di orbit planet Mars siang ini?" \textit{(Pertanyaan bukan mengklaim fakta, ia menagih jawaban. Tak bernilai T/F.)}
    \item \textbf{Kalimat Imperatif (Perintah/Larangan):} "Segera simpan \textit{file} dokumen itu, atau Anda akan kena marah Dosen!" \textit{(Ini murni perintah, tak memiliki kepastian nilai matematis.)}
    \item \textbf{Kalimat Seruan (Eksklamatif):} "Aduhai, alangkah cantiknya mahakarya patung lilin museum tersebut!" \textit{(Frasa seruan subjektif keindahan tidak bisa dipatok universal False/True.)}
\end{enumerate}
\end{contoh}

\subsection{Kalimat Terbuka (Open Sentences)}

Terdapat satu garis batas mematikan yang patut diwaspadai: persimpangan antara Kalimat Deklaratif dengan \textbf{Kalimat Terbuka}. 

\begin{definisi}[Kalimat Terbuka]
Adalah rangkaian sintaks narasi aljabar yang \textbf{MENGANDUNG VARIABEL BEBAS (Peubah)}, sehingga sistem otak kita belum sanggup menghakimi apakah klaimnya itu sah True atau False sampai nyawa si peubah itu kita ganti substitusi angka definitif.
\end{definisi}

\begin{contoh}[Misteri Peubah Rahasia]
    \begin{itemize}
        \item "$x + 7 = 12$" \quad $\rightarrow$ Bukan bagian dari Proposisi! Kita lumpuh menimbang nilainya. Kalau dewa nasib melempar $x = 5$, si dia mekar jadi BENAR. Kalau si dewa iseng menukar input merusak jadi $x=10$, dia mendadak cacat menjadi SALAH.
        \item "$He$ adalah mahasiswa teladan paling cerdas se-antero kelas Logika." \quad $\rightarrow$ Frasa "$He$" adalah sapaan anonim tak bertuan! Entah dia merujuk pada Budi (Mungkin True) atau malahan si pemalas Joni (Murni False). Statusnya haram diangkat menjadi proposisi.
    \end{itemize}
\end{contoh}
Mematangkan insting menyortir kalimat deklaratif bermakna tegas (\textbf{Proposisi Tertutup}) dari gunungan *Non-Proposisi* ganda makna merupakan insting perdana seorang programer meracik spesifikasi algoritma perangkat keras tanpa cela ambiguitas.
